\section{Results and Discussions}\label{sec:results}

Using the RITDs extracted in the previous section, we examine the pion-mass and lattice-spacing dependence.
The top of Fig.~\ref{fig:RITD-ensembles} shows the $\eta_s$ RITDs at boost momentum around 2~GeV as functions of the Wilson-line length $z$ for the a12m220, a12m310, and a15m310 ensembles.
We see no noticeable lattice-spacing dependence.
The bottom of Fig.~\ref{fig:RITD-ensembles} shows the pion RITDs with boost momentum around 1.3~GeV for the same ensembles.
Again, there is no visible lattice-spacing or pion-mass dependence.


\begin{figure}[htbp]
\centering
	\centering
	\includegraphics[width=0.45\textwidth]{images/ITD_etas_p456_comparision.pdf}
	\includegraphics[width=0.45\textwidth]{images/ITD_pion_p334_comparision.pdf}
\caption{The $\eta_s$ (top) and pion (bottom) RITDs at boost momenta $P_z \approx 2$~GeV and 1.3~GeV, respectively, for the a12m220, a12m310, and a15m310 ensembles.
In both cases, we observe weak lattice-spacing and pion-mass dependence.
}
\label{fig:RITD-ensembles}
\end{figure}


To extract gluon PDFs, we follow the steps in Sec.~\ref{sec:cal-details} between Eq.~\ref{eq:ME_unpol} and Eq.~\ref{evolution} by first obtaining EITDs and using Eq.~\ref{conversion} to extract $g(x)$.
To obtain EITDs, we need the RITD $\mathscr{M}(\nu,z^2)$ to be a continuous function of $\nu$ to evaluate the $x\in[0,1]$ integral in Eq.~\ref{evolution}.
We achieve this by using a ``$z$-expansion''\footnote{Note that the $z$ in the ``$z$-expansion'' is not related to the Wilson link length $z$ we use elsewhere.} fit~\cite{Boyd:1994tt,Bourrely:2008za} to the RITD.
The following form is used \cite{Joo:2019bzr}:
%
\begin{equation}
\mathscr{M}(\nu, z^2,M_{\pi}) = \sum_{k=0}^{k_\text{max}}\lambda_k\tau^k,
\label{zexpansion}
\end{equation}
%
where $\tau=\frac{\sqrt{\nu_\text{cut}+\nu}-\sqrt{\nu_\text{cut}}}{\sqrt{\nu_\text{cut}+\nu}+\sqrt{\nu_\text{cut}}}$.
Then, we use the fitted $\mathscr{M}(\nu, z^2)$ in the integral in Eq.~\ref{evolution}.
The $z$-dependence in the $\mathscr{M}(u\nu, z^2)$ term of the evolution function comes from the one-loop matching term, which is a higher-order correction compared to the tree-level term;
thus, the $z$-dependence can be neglected in $\mathscr{M}(\nu, z^2)$.
We adopt as the best value $\nu_\text{cut}=1$, as used in Ref.~\cite{Joo:2019bzr}, but we also vary $\nu_\text{cut}$ in the range $[0.5,2]$, and the results are consistent.
We fix $\lambda_0 = 1$ to enforce the RITD $\mathscr{M}(\nu, z^2)$ in Eq.~\ref{eq:RITD}.
The expansion order $k_\text{max}=3$ is used, because we can fit all the data points of $P_z\in [1,5]\times 2\pi/L$ ($P_z\in [1,7]\times 2\pi/L$ for a12m220 ensemble)
and $z$ up to 0.6~fm
with $\chi^2/\text{dof}<1$ using a 4-term $z$-expansion for each ensemble.
The reconstructed bands from ``$z$-expansion'' on RITDs are shown in the upper plot in Fig.~\ref{fig:ITD-comparison}.
They describe the RITD data points well for all ensembles.


\begin{figure}[htbp]
\centering
	\centering
	\includegraphics[width=0.45\textwidth]{images/RITD_pion_comparision.pdf}
	\includegraphics[width=0.45\textwidth]{images/EITD_pion_comparision.pdf}
\caption{
The RITDs $\mathscr{M}$ with reconstructed bands from ``$z$-expansion'' fits (top) and the EITDs $G$ with reconstructed bands from fits (bottom) calculated on ensembles with lattice spacing $a\approx 0.12$~fm, pion masses $M_\pi\approx\{220,310,690\}$~MeV, and $a\approx 0.15$~fm, $M_\pi\approx 310$~MeV, noticing that $a\approx 0.12$, $M_\pi\approx 690$~MeV results are from a12m220 ensemble here.
}
\label{fig:ITD-comparison}
\end{figure}


After we have the continuous-$\nu$ fitted RITDs, we obtain the EITDs through Eq.~\ref{evolution}.
The RITDs $\mathscr{M}$ and EITDs $G$ as functions of $\nu$ on all ensembles studied in this work are shown in Fig.~\ref{fig:ITD-comparison}.
At some $\nu$ values, there are multiple $z$ and $P_z$ combinations for a fixed $\nu$ value.
Therefore, there are points in the same color and symbol overlapping at the same $\nu$ from the same lattice spacing and pion mass.
To match with the lightcone gluon PDF through Eq.~\ref{conversion}, the EITDs $G(\nu,\mu)$ should be free of $z^2$ dependence.
However, the EITDs obtained from Eq.~\ref{evolution} have $z^2$ dependence from neglecting the gluon-in-quark contribution and higher-order terms in the matching.
The EITDs also depend on lattice-spacing $a$ and pion-mass $M_\pi$.
Recall that the RITDs show weak dependence on lattice spacing $a$ and pion mass $M_\pi$.
We see that the effects of $a$ and $M_\pi$ dependence on the EITDs are also not large;
the EITD results from different $a$, $M_\pi$ are mostly consistent with each other, as shown in the second row of Fig.~\ref{fig:ITD-comparison}.
We also observe a weak dependence on $z^2$ for the RITDs and EITDs in Fig.~\ref{fig:ITD-comparison}.

The gluon PDF $g(x,\mu^2)$ can now be extracted from the EITDs using Eq.~\ref{conversion}.
%
We assume a functional form, also used by JAM~\cite{Barry:2018ort,Cao:2021aci}, for the lightcone PDF to fit the EITD,
%
\begin{align}
f_g(x,\mu) = \frac{xg(x, \mu)}{\langle x \rangle_g(\mu)} = \frac{x^A(1-x)^C}{B(A+1,C+1)},
\label{functional}
\end{align}
%
for $x\in[0,1]$ and zero elsewhere.
The beta function $B(A+1,C+1)=\int_0^1 dx\, x^A(1-x)^C$ is used to normalize the area to unity.
Then, we apply the matching formula to obtain the EITD ${G}$ from the functional form PDF using Eq.~\ref{conversion}.
We fit the EITDs $G(\nu,\mu)$ obtained from the parametrization to the EITDs $G(\nu,z^2,\mu,a,M_\pi)$ from the lattice calculation.


\begin{figure}[htbp]
\centering
\centering
\includegraphics[width=0.45\textwidth]{images/xg-x0-1_pion_a12m220_par.pdf}
\caption{
The $xg(x, \mu)/\langle x \rangle_g$ at $\mu^2=4\text{ GeV}^2$ as function of $x$ (bottom) calculated with lattice spacing $a\approx 0.12$~fm, pion masses $M_\pi\approx 220$~MeV with the fitted bands of  $z_\text{max}\approx 0.6$~fm from the 1-, 2- and 3-parameter fits described in Eq.~\ref{functional} and the paragraph after it.
}
\label{fig:fitform-comparison}
\end{figure}


We investigate the systematic uncertainty introduced by the different parametrization forms which are commonly used for $f_g(x,\mu)$ in PDF global analysis and some lattice calculations.
The first one is the 2-parameter form in Eq.~\ref{functional}.
Second, we consider the 1-parameter form $N_1 (1-x)^C$ used in xFitter's analysis~\cite{Novikov:2020snp} (also used in Ref.~\cite{Aurenche:1989sx,Sutton:1991ay}), which is equivalent to Eq.~\ref{functional} with $A=0$.
Third, we consider a 3-parameter form $N_3 x^A(1-x)^C(1+D\sqrt{x})$.
We fit the three different forms to the EITDs of lattice data with $z_\text{max}\approx 0.6$~fm by applying the scheme conversion Eq.~\ref{conversion} to the 1-, 2- and 3-parameter PDF forms.
Here, we focus on the result from the lightest pion mass $M_{\pi}\approx 220$~MeV at lattice spacing $a\approx 0.12$~fm.
The $\chi^2/\text{dof}$ of the fits decreases as $1.47(72)$, $1.08(68)$, to $1.04(41)$, shows slightly better fit quality for 2- and 3-parameter fits.
As shown in Fig.~\ref{fig:fitform-comparison}, there is a big discrepancy between the $f_g(x,\mu)$ fit bands from the 1-parameter fit and the 2-parameter fit in the $x<0.4$ region, but the discrepancy between the 2- and 3-parameter fits is much smaller.
Therefore, we conclude that 1-parameter fit on lattice data here is not quite reliable, and the fit results converge at the 2- and 3-parameter fits.
The same conclusions hold for all other ensembles and pion masses.
Therefore, using the 2-parameter form defined Eq.~\ref{functional} (same parametrization as JAM) for our final results is very reasonable.


Another source of systematic uncertainty comes from neglecting the contribution of the quark term in Eq.~\ref{matching-eq} based on the assumption (motivated by global fits) that the pion $q_S(x)$ is smaller than the gluon PDF.
Currently, there are no $q_S(x)$ results from lattice simulation since only the valence distribution of the pion has been done.
Thus, we estimate the systematic due to omitting the $q_S(x)$ contribution by using the pion quark PDFs from xFitter~\cite{Novikov:2020snp} at NLO.
Using these, we obtain revised RITDs and EITDs including the gluon-in-quark $R_{gq}$ term focusing on example from the $a\approx 0.12$~fm, pion mass $M_\pi\approx 220$~MeV lattice, repeating the same procedure from Eq.~\ref{zexpansion} and fitting the EITDs with Eq.~\ref{functional}.
On the right-hand side of Fig.~\ref{fig:xgx}, we show the mean value of $xg(x, \mu)/\langle x \rangle_g$ with both gluon-in-gluon (gg) and gluon-in-quark (gq) contributions (the blue solid line) compared to the a12m220 results using the gluon-in-gluon contribution only (the green solid line).
There are 5 to $10\%$ differences in the mean value including the gluon-in-gluon contribution for $x<0.9$, which indicates that the gluon-in-quark contribution is relatively small at $\mu^2=4\text{ GeV}^2$ compared to the current statistical errors in the small-$x$ region.
In the $x>0.9$ region, the gluon-in-quark contribution becomes more significant, but it remains smaller than the statistical error.
Once studies are available with sufficiently reduced statistical uncertainty in the large-$x$ region, the quark contribution will need to be included.

From the above analyses of the choice of fit form and the contribution of the quark term, we conclude that these systematics are negligible relative to the current statistics.
Therefore, we adopt the $z_\text{max}\approx 0.6$~fm ($z_\text{max}\approx 0.75$~fm for a15m310 ensembles) fits to the EITDs, neglect the quark contribution term in the matching, and use the Eq.~\ref{conversion} fit form for our final results on all lattice ensembles.
The $xg(x, \mu)/\langle x \rangle_g$ reconstructed fit bands of these ensembles are shown in the left plot in Fig.~\ref{fig:xgx}, comparing results from different lattice spacings and pion masses.
The reconstructed fit bands with different pion mass $M_\pi\approx\{220,310,690\}$~MeV are consistent at the same lattice spacing $a\approx 0.12$~fm, indicating mild gluon PDF dependence on pion mass.
Similarly, when comparing lattice-spacing dependence of pion PDFs using data around pion mass $M_\pi\approx 310$~MeV, we find that fitted PDF is slightly smaller in the $x>0.1$ region for the 0.12-fm lattice, but still within one sigma, which indicates the lattice-spacing dependence is also mild.
We also note that the bands from different ensembles show a differing speed of fall-off as $x\to 1$ in the large-$x$ region.
We study this fall-off behavior in more depth below.

The behavior of the gluon PDF fall-off in the large-$x$ region is widely studied in both theory and global analyses.
Perturbative QCD studies~\cite{Close:1977qx,Sivers:1982wk} and DSE calculations~\cite{Bednar:2018mtf,Freese:2021zne} suggest that the gluon distribution $g(x,\mu^2) \sim (1-x)^C$ with $C\approx 3$ in the limit $x\to 1$.
The prediction from perturbative QCD~\cite{Sivers:1982wk} is based on the idea that the gluon PDF should be suppressed at large $x$ relative to the quark PDF, because the quarks are the sources of large-$x$ gluons;
that is, $g(x,\mu^2)/q_v(x,\mu^2)\to 0$ as $x\to 1$.
Early fits of experimental data gave $C\approx 2$~\cite{Aurenche:1989sx,Sutton:1991ay} or $C<2$~\cite{Gluck:1991ey,Gluck:1999xe}, but the more recent global analysis from JAM collaboration yielded $C>3$~\cite{Barry:2018ort,Cao:2021aci} and xFitter collaboration found $C\approx 3$~\cite{Novikov:2020snp}.
Our fitted parameter $C$ is $3.6(1.5)$, $3.3(2.0)$, $4.7(2.8)$ for $M_\pi\approx\{690,310,220\}$~MeV, respectively, at lattice spacing $a\approx 0.12$~fm.
These $C$ results are consistent with each other and show a slightly increasing trend as the pion mass approaches the physical pion mass.
For lattice spacings $a\approx\{0.15,0.12\}$~fm, $C=\{2.2(1.5),3.3(2.0)\}$, respectively, at $M_\pi\approx 310$~MeV, which suggests that $C$ will increase toward the continuum limit.
We also investigate the effect of the gluon-in-quark contribution on the $C$ value, and it makes about $0.1$ difference, which we neglect.
Given that both the pion-mass and lattice-spacing extrapolations seem to show increasing $C$, it seems reasonable to conclude from this lattice-QCD study that $C>3$.

We compare our reconstructed gluon PDF to those from global fits on the right-hand side of Fig.~\ref{fig:xgx}.
It shows the $xg(x, \mu)/\langle x \rangle_g$ reconstructed fit band of $a\approx 0.12$~fm, $M_\pi\approx 220$~MeV lattice and the NLO pion gluon PDFs from xFitter~\cite{Novikov:2020snp} and JAM~\cite{Barry:2018ort,Cao:2021aci} at $\mu^2=4\text{ GeV}^2$.
The JAM band appears somewhat wider than expected, because we reconstruct it by dividing $xg(x, \mu)$ by the mean value of $\langle x \rangle_g$;
the correlated values needed for a correct error estimation were not available.
Note that xFitter uses the fit form of Eq.~\ref{functional} with $A=0$.
Our fitted pion gluon PDF is consistent with JAM for $x>0.2$ and with xFitter for $x>0.5$ within one sigma.
We see in this comparison that our results are of similar error size as the  global-fit analysis and are useful to provide constraints from theoretical calculation in addition to the experimental data.


\begin{figure*}[htbp]
\centering
	\centering
	\includegraphics[width=0.45\textwidth]{images/xg-x0-1_bandcomp_aM_pion.pdf}
	\includegraphics[width=0.45\textwidth]{images/xg-x0-1_bandcomp_JX_pion.pdf}
\caption{
The pion gluon PDF $xg(x, \mu)/\langle x \rangle_g$ as a function of $x$ obtained from the fit to the lattice data on ensembles with lattice spacing $a\approx\{0.12,0.15\}$~fm, pion masses $M_\pi\approx\{220,310,690\}$~MeV (left), and $xg(x, \mu)/\langle x \rangle_g$ as function of $x$ obtained from lattices of $a\approx 0.12$~fm, $M_\pi\approx 220$~MeV (right), compared with the NLO pion gluon PDFs from xFitter and JAM at $\mu=2$~GeV in the $\overline{\text{MS}}$ scheme.
The JAM error shown is overestimated due to lack of available correlated uncertainties in its constituent components.
Our PDF results are consistent with JAM~\cite{Barry:2018ort,Cao:2021aci} for $x>0.2$ and xFitter~\cite{Novikov:2020snp} for $x>0.5$.
}
\label{fig:xgx}
\end{figure*}
